\section{Recommendation}
\label{sec:recommendation}

Proceed --- in the reshaped form this study argues for, and in the
phased order that front-loads standalone value.

\paragraph{Adopt.} An embedded, tactic-oriented proof-script DSL in
the lineage of \code{dotnetdll}'s \code{asm!}~\cite{dotnetdll}:
block-shaped \code{ves::proof!} scripts wrapping each \code{unsafe}
site, premises acquired as zero-sized ghost tokens whose existence,
\code{cfg}-liveness, and dataflow rustc itself checks, goals drawn
from (or inferred against) a single declarative vocabulary that
generates both the Rust tokens and the Lean declarations. Scripts
elaborate deterministically to Lean~4 statements; proofs are ordinary
Lean terms driven by reusable domain tactics that are untrusted
metaprograms; the kernel is Lean's; Curry--Howard is inherited, not
reimplemented. The whole verifier lives in a separately versioned
repository whose only contact with \dotnetrs{} is two tiny Rust
dependencies and two hashed manifest artifacts
(\cref{sec:design-compart}) --- the VM's build never touches Lean.
Surface conventions stay aligned with RFC 3842 safety
tags~\cite{rfc3842} so the project tracks, and can eventually
contribute to, the community convention.

\paragraph{Reject.} A bespoke proof kernel and solver
(\cref{sec:design-custom}: unauditable TCB, automation cost dwarfs the
kernel, forfeits the agent ecosystem) --- including its subtle
reappearance as a ``smart'' checker that evaluates tactics outside
Lean; wholesale adoption of any existing verifier as the whole answer
(\cref{sec:design-adopt}: none reaches the ECMA-335 layer); attribute
or comment syntax as the primary carrier for multi-step proofs (kept
only as the degenerate single-lemma form); and any plan whose first
deliverable is a proof rather than the token layer and the model
(\cref{sec:risks-tcb}: untethered proofs decay into documentation).

\paragraph{Immediate actions, independent of any spin-off decision.}
Phase 0 is due diligence, not verification: fix the two live defects,
promote at least one Miri leg to blocking, and expand the differential
harness past its single test. Phase 1 (the embedded layer) is the
cheapest experiment that tests the whole thesis on real ground ---
and, uniquely among the phases, it yields value even with no Lean
involvement: ghost tokens make premise citations compile-checked in
every feature configuration, which would already have made defect (a)
unwritable.

\paragraph{Two genuinely open decisions.} First, how far to trust
the $\mathcal{M}_{\mathrm{Rust}}$ axiomatization before reaching for
Rocq foundations --- this study says: cap it at $\sim$30
Miri-falsifiable axioms and revisit after the critical sites are proved.
Second, whether the VES model targets the standard's semantics or
\dotnetrs{}'s documented deviations (single load context, non-moving
GC, STW-only collection) --- this study says: model the standard, and
register each deviation as a named trust grant, because the deviations
are exactly the things most likely to change and the grant registry is
what makes changing them safe.

\medskip

\noindent The strongest reason to proceed is internal to the
project. \dotnetrs{} has pushed Rust-native assurance --- branded
lifetimes, capability tokens, compile-time lock DAGs, validation
magics, a written proof charter --- about as far as it goes, and the
residue is a well-defined, finite set of obligations that are theorems
or they are nothing. The project's own documentation already writes
them as theorems in prose (``STW keeps every arena alive and every
mutator stopped through this read''); its own metadata layer already
demonstrates the embedding idiom the proof language should use. The
proposal, correctly scoped, is not adding a formal layer to an
informal project; it is giving a project that already reasons formally
a checker for the reasoning it already does --- and producing, along
the way, the first mechanized model of the specification it
implements.
